\documentclass[../../main.tex]{subfiles}
\begin{document}

% 年度の大きな表示
\begin{center}
  {\fontsize{48pt}{58pt}\selectfont \textbf{1994}\normalsize\textbf{年度}}
\end{center}

\vspace{10mm}

% 本文


\vspace{15mm}

% 出題テーマと難易度の表
\noindent\textbf{出題テーマと難易度}

\vspace{5mm}

\begin{center}
  \shadowbox{%
    \begin{tabular}{|c|p{8cm}|c|}
      \hline
      \textbf{問題番号} & \textbf{テーマ} & \textbf{難易度} \\
      \hline
      第1問 & 放物線と弦で囲まれる面積一定と接線交点の軌跡 & \\
      \hline
      第2問 & 双曲線上の点の回転と軌跡 & \\
      \hline
      第3問 & 指数関数と三角関数の積の定積分・無限級数 & \\
      \hline
      第4問 & 整数の組に対する関数の値と一意性の証明 & \\
      \hline
      第5問 & 2つの放物線上の動点間の距離の最小値 & \\
      \hline
    \end{tabular}%
  }
\end{center}

\vspace{15mm}

% 解答
\noindent\textbf{解答}

\vspace{5mm}

\begin{center}
  \shadowbox{%
    \renewcommand{\arraystretch}{1.6}
    \begin{tabular}{|c|c|p{8cm}|}
      \hline
      \textbf{問題番号} & \textbf{小問} & \textbf{答え} \\
      \hline
      第1問 & - & $y=x^2-\dfrac{1}{4}$ \\
      \hline
      \multirow{3}{*}{第2問} & (1) & $Q\left(t\cos\theta+\dfrac{2}{t}\sin\theta,\ t\sin\theta-\dfrac{2}{t}\cos\theta\right)$ \\
      \cline{2-3}
                             & (2) & $(x\cos\theta+y\sin\theta)(-x\sin\theta+y\cos\theta)=-2$ \\
      \cline{2-3}
                             & (3) & $\theta=\dfrac{\pi}{6},\dfrac{\pi}{2},\dfrac{7\pi}{6},\dfrac{3\pi}{2}$ \\
      \hline
      \multirow{2}{*}{第3問} & (1) & $\dfrac{1}{2}(1+e^{-\pi})$ \\
      \cline{2-3}
                             & (2) & $\dfrac{1}{2}\cdot\dfrac{1+e^{-\pi}}{1-e^{-\pi}}$ \\
      \hline
      \multirow{2}{*}{第4問} & (1) & $(m,n)=(0,0),(0,1),(0,2),(1,0),(1,1),(2,0)$ \\
      \cline{2-3}
                             & (2) & 証明問題 \\
      \hline
      第5問 & - & $\min|PQ|=3\sqrt{5}$ \\
      \hline
    \end{tabular}%
    \renewcommand{\arraystretch}{1.0}
  }
\end{center}

\vspace{15mm}

% 解答時間
\noindent\textbf{試験形態}

\vspace{5mm}

\begin{center}
  \shadowbox{%
    \begin{tabular}{p{8cm}}
      （要確認）
    \end{tabular}%
  }
\end{center}

\end{document}
